% Complete proof candidate. Individual check status is in result_audit.json.
\subsection{C.6. Supplement: ranks, saturation, and central derivatives}

We retain the order of the eleven rational CM values in the prompt, and
write $E_i$ for the global minimal models recorded in the preceding
tables. The real period in this subsection is
\[
 \Omega_i=\int_{E_i(\R)}|\omega_i|,
\]
including both connected components when $E_i(\R)$ is disconnected.
The canonical height $\widehat h$ is relative to $(O)$; PARI's
\texttt{ellheight} and the regulator convention used in the BSD formula
are relative to $2(O)$, and hence equal $2\widehat h$ in rank one.
This factor of two is essential in every comparison below.

\begin{proposition}[ARI C.6.R1; Proved by descent and full saturation]\label{ari:C.6.R1}
The algebraic and analytic ranks agree for the eleven curves. They are
two for $j=8000$ and $j=-12288000$, and one for all nine other values.
For every rank-one curve the section $P$ is a primitive generator of
$E(\Q)/E(\Q)_{\rm tors}$. Its index is therefore one. For the two
rank-two curves the index of $\Z P$ in this free quotient is infinite.
\[
\begin{array}{r|r|r|r}
i&j&\operatorname{rank}E(\Q)&[E(\Q)/E(\Q)_{\rm tors}:\Z P]\\\hline
1&8000&2&\infty\\
2&54000&1&1\\
3&287496&1&1\\
4&16581375&1&1\\
5&-3375&1&1\\
6&-32768&1&1\\
7&-884736&1&1\\
8&-12288000&2&\infty\\
9&-884736000&1&1\\
10&-147197952000&1&1\\
11&-262537412640768000&1&1
\end{array}
\]
\end{proposition}
\begin{proof}[Arithmetic verification]
Exact $2$-descent in PARI/GP returns equal lower and upper rank bounds
for all eleven curves; its entire output, including independent points,
is in \path{code/C_cm_rank_lvalues.out}. Independent eclib descent
and saturation were completed for the first nine curves. The separate
full-saturation calculation \path{code/Crank_saturation.cpp} treats
all nine rank-one points, using eclib's automatic height-based index
bound and checking every prime up to that bound. All nine returned
index one and an empty list of unsaturated primes.

For clarity, the rounded-down bounds for the indices in the subgroup
of points having good reduction everywhere are respectively
\[
 1,1,1,2,2,2,3,5,5
\]
in the rank-one order $i=2,3,4,5,6,7,9,10,11$; the algorithm also
checks the primes dividing component-group orders. The underlying
height lower bound is the one in
\citet[Theorem 1]{CrCremonaSiksek2006}, applied to a minimal
Weierstrass equation and the subgroup
$E(\Q)\cap E^0(\R)\cap\bigcap_pE^0(\Q_p)$.
For the largest curve the lower bound on the nonzero height in this
subgroup is $0.651940632595967949\ldots$, its regulator is
$17.323913974368488616\ldots$, and the resulting index bound is
$5.154885040175890631\ldots$. The finite-field saturation maps for
$2,3,5$ already have full rank at residue primes $17,17,59$.
These data certify primitivity; merely finding $P$ in the output of
\texttt{ellrank} would not do so. The complete height-bound and
saturation transcripts are retained, including the independent
descent run that was stopped after its ninth completed curve.

The analytic ranks follow from the explicit derivative sums and bounds
proved below. The two rank-two bases furnished by descent are
$(-7,441),(105,833)$ for $i=1$, and
$(0,32004),(22218,3305456)$ for $i=8$.
\end{proof}

\begin{proposition}[ARI C.6.R2; Proved: central derivatives with explicit tails]\label{ari:C.6.R2}
The following are the values of $L'(E_i,1)/\Omega_i$; zeros in rows
one and eight are exact. More digits, the periods, derivatives and
tail bounds are in \path{results_Crank.csv}.
\begin{center}\small
\begin{tabular}{r|l}
$i$&$L'(E_i,1)/\Omega_i$\\\hline
1&$0$\\
2&$5.134398425737889664345619204085016979705061236171249722654226\ldots$\\
3&$2.134120849044754020051650801945664635106456422166881627967598\ldots$\\
4&$31.81013979419927776012209427941545866492237285892013668350005\ldots$\\
5&$1.869462708725885120030080805051387134542481374813232701178304\ldots$\\
6&$3.428873379050567976208155749688133688667614469824389504923957\ldots$\\
7&$0.362583348071598548315763082767526599722077924892872839203426\ldots$\\
8&$0$\\
9&$9.629273918234131678053742152263035417231796757756973886621556\ldots$\\
10&$16.04238857830583792513538634637663343755958963581574670348088\ldots$\\
11&$1247.321806154531180413038860827393095408344267114181914400926\ldots$
\end{tabular}
\end{center}
The two second derivatives, divided by the same real periods, are
\begin{align*}
 L''(E_1,1)/\Omega_1
 &=28.71408722758099118128474692322079946598060512912402172882525\ldots,\\
 L''(E_8,1)/\Omega_8
 &=78.03222041637032223732443782051140882655135939512581188900823\ldots.
\end{align*}
In particular the leading Taylor coefficients for these two curves
are half the displayed second derivatives.
\end{proposition}
\begin{proof}[Derivative formulas and error control]
Let $N$ be the conductor, $c=2\pi/\sqrt N$, and
$L(E,s)=\sum_{n\geq1}a_nn^{-s}$. Put
\[
 E_1(x)=\int_x^\infty\frac{e^{-u}}u\,du,
 \qquad G(x)=\int_x^\infty e^{-u}\frac{\log(u/x)}u\,du.
\]
Splitting the Mellin integral at $1/\sqrt N$ and using the functional
equation gives, for root number $-1$,
\begin{equation}\label{ari:eq:Cr-first}
 L'(E,1)=2\sum_{n\geq1}\frac{a_n}{n}E_1(cn).
\end{equation}
For root number $+1$ and $L(E,1)=0$, it gives
\begin{equation}\label{ari:eq:Cr-second}
 L''(E,1)=4\sum_{n\geq1}\frac{a_n}{n}G(cn).
\end{equation}
One can verify the factors directly by writing the completed Mellin
integral as an integral over $[1,\infty)$ of
$(u^{s-1}+wu^{1-s})f(iu/\sqrt N)$: its first derivative has
the factor $2\log u$ when $w=-1$, and its second derivative
the factor $2\log^2u$ when $w=1$. Integration by parts gives
the two formulas above. The vanishing of the lower derivatives
removes all derivatives of the gamma and conductor factors.

Hasse's bound and multiplicativity imply
$|a_n|\leq d(n)\sqrt n\leq2n$. Also
$E_1(x)\leq e^{-x}/x$ and $G(x)\leq e^{-x}/x^2$.
The latter follows on writing $u=x+v$ and using
$\log(1+v/x)\leq v/x$. Consequently the absolute tails after
$n=M$ are at most
\begin{equation}\label{ari:eq:Cr-tails}
 \frac{4e^{-c(M+1)}}{c(M+1)(1-e^{-c})},\qquad
 \frac{8e^{-c(M+1)}}{c^2(M+1)^2(1-e^{-c})},
\end{equation}
respectively. For the first ten curves we used
$M=\lceil180/c\rceil$, exact integer coefficients, and 1024-bit
MPFR arithmetic. For the rank-two kernel we used the convergent series
\[
 G(x)=\frac12(\log x+\gamma)^2+\frac{\pi^2}{12}
       +\sum_{k\geq1}\frac{(-x)^k}{k^2k!}.
\]
After $k>x$ its terms decrease in absolute value. The alternating
remainder was cut below $10^{-140}$. At most 740 terms were needed;
the largest intermediate term was less than $1.4\cdot10^{72}$.
The 1024-bit arithmetic therefore leaves far more than 100 digits
after this cancellation. The largest bound in
\eqref{ari:eq:Cr-tails} was $3.46\cdot10^{-77}$. The source and complete
outputs are \path{code/Crank_coefficients.gp},
\path{code/Crank_derivatives.cpp}, and their corresponding outputs.

For $i=11$, a full vector of $819$ million coefficients would be
unnecessarily large. We instead evaluated the CM theta series in
bounded-memory segments. The curve is the quadratic twist, by
$d=185801=7\cdot11\cdot19\cdot127$, of the conductor-$163^2$
curve
\[
 y^2+y=x^3-2174420x+1234136692.
\]
Write $D=163$, $\chi_D(a)=(a/D)$, and $\chi_d(n)=(d/n)$.
The exact coefficient $a_n$ is the sum of
\[
 \chi_d(n)\chi_D(a)a \quad(n=a^2,\ a>0)
\]
and
\[
 \chi_d(n)\chi_D((D+1)u/2)u
 \quad\left(n=\frac{u^2+Dv^2}{4},\ u,v>0,
                           \ u\equiv v\pmod2\right).
\]
This is the Hecke-character theta series for
$\mathcal O_{\Q(\sqrt{-D})}$, grouping conjugate algebraic integers
and then applying the quadratic twist; see
\citet[Section 3]{CrDarmon2009}. The first 10000 coefficients agree
exactly with PARI for both the untwisted and twisted curve. The same
code at $D=67,d=217$ reproduces the independent PARI derivative.

To accelerate the sum, in a block $n=n_0+r$, $|r|\leq R$, form the
exact integer moments $M_k=\sum a_nr^k$. For $x=cn_0$ write
\[
 e_k=e^{-x}(-x)^k/k!,\quad g_0=e^{-x},\quad
 g_k=e_k-g_{k-1},\quad
 \phi_0=E_1(x),\quad\phi_k=-\phi_{k-1}-g_{k-1}/k.
\]
Then the block contribution to $L'/2$ is
$\sum_k\phi_k M_k/n_0^{k+1}$. This follows by Taylor expansion
of $E_1(x(1+s))/(1+s)$. On $|s|=1/2$ its modulus is at most
$2E_1(x/2)$, which is less than $14$ throughout the block range.
We take enough terms that $(2R/n_0)^{k+1}<10^{-85}$.
The Cauchy remainder, summed with the actual
$\sum|a_n|<2\cdot10^{12}$ and $n_0>7700$, is less than
$10^{-75}$. The values of $E_1$ at consecutive block centres
are continued by the analogous Taylor expansion and refreshed
from a correctly rounded MPFR evaluation every 1024 blocks.
The continuation bound is $10^{-110}$ per normalized expansion.
All moments are integers, not floating-point sums.

The 384-bit run used $R=128$ and cutoff $170$; a separate 512-bit
run used $R=96$ and cutoff $190$. Their values agree beyond 75
decimal places. The latter used $915818441$ coefficients in
segments, with an omitted tail less than $3.10\cdot10^{-78}$,
and gave
\[
 L'(E_{11},1)=
 2.296582198163380608965915774436446055185027904705664673866754\ldots.
\]
The elementary operation-rounding bounds at these precisions are
smaller than $10^{-90}$: in the recurrences above
$|e_k|\leq1$, $|g_k|\leq1$, and $|\phi_k|<20$;
there are fewer than $5\cdot10^6$ blocks and the recurrences are
refreshed before 1024 consecutive continuations. Bounding each
normalized moment by $\sum|a_n|(R/n_0)^k$ gives this estimate
without assuming cancellation in the total sum. The arithmetic
and truncation margins are amply sufficient for every displayed
50-digit value. Both executable builds and both transcripts are
retained as \path{code/Crank_cm_stream*}.

The exact root numbers are $+1$ for $i=1,8$ and $-1$ otherwise.
All points are nontorsion by Statement C.5. For a CM curve,
the Coates--Wiles theorem, in the form
\citet[Theorem 10.1]{CrRubin1999}, says that a nonzero central
Hecke $L$-value forces the Mordell--Weil group over the CM field
to be finite. For a nonmaximal CM order first pass to an isogenous
curve with maximal order; ranks and $L$-functions are unchanged.
Since $L(E/K,s)=L(E/\Q,s)^2$ for the CM field $K$, nontorsion
therefore forces $L(E,1)=0$. This proves the needed central
vanishing also in the two sign-$+1$ cases. The functional equation
then forces $L'(E,1)=0$ for those two curves. The positive first
or second derivatives just evaluated prove exactly the asserted
analytic ranks.
\end{proof}

\begin{proposition}[ARI C.6.R3; Proved: the consequences of Gross--Zagier and Kolyvagin]\label{ari:C.6.R3}
For each of the nine rank-one curves there is a positive rational
number $q_i$ such that
\begin{equation}\label{ari:eq:Cr-rational-factor}
 \frac{L'(E_i,1)}{\Omega_i}=q_i\widehat h(P_i).
\end{equation}
Their Tate--Shafarevich groups over $\Q$ are finite. Consequently
the corresponding $D_{\rm part}$ is an explicit logarithm minus
$(6/(\pi q_i))L'(E_i,1)/\Omega_i$. The two rank-two curves do
not satisfy a rank-one formula; there $L'(E_i,1)=0$ while
$\widehat h(P_i)>0$.
\end{proposition}
\begin{proof}
Gross--Zagier's theorem applies because each curve is modular and
$L(E_i,1)=0$. CM modularity follows from its Hecke-character
theta series. Their Theorem I.7.3, in the opened original
preprint \citep{CrGrossZagier1986}, gives a rational point $R$ and
a nonzero rational factor relating $L'(E,1)$ to its canonical
height times the real period. Changing their period convention
to the full real period changes that factor only rationally.
As $L'(E_i,1)>0$, $R$ has infinite order. Primitivity gives
$R=mP_i+T$, $m\ne0$, $T$ torsion; its height is
$m^2\widehat h(P_i)$. This proves
\eqref{ari:eq:Cr-rational-factor}, including positivity.

The rank-one Gross--Zagier--Kolyvagin implication asserts rank one
and finite Tate--Shafarevich group when the central zero is simple.
One can use Kolyvagin's Heegner-point theorem as stated in
\citet[Theorem 1.3]{CrGross1991}, together with the nonzero
Gross--Zagier Heegner height; a modern primary research statement
of the rank conclusion is \citet[Theorem 1.1]{CrJetchev2008}.
The hypotheses are satisfied separately by the positive derivative
and exact sign in each of the nine rows. The two remaining rows
have analytic rank two, and this implication makes no assertion
of a height formula or of finiteness of their Tate--Shafarevich
groups. In particular it cannot express the second derivative
in terms of the height of $P$ alone: the rank-two regulator involves
a second independent point and their height pairing.
\end{proof}

\paragraph{What the numerical BSD comparison does and does not determine.}
For a primitive rank-one point, the BSD prediction in these conventions is
\[
 \frac{L'(E,1)}{\Omega_E}
 =\frac{2\#\operatorname{Sha}(E/\Q)\prod_pc_p}
        {\#E(\Q)_{\rm tors}^{\,2}}\,\widehat h(P).
\]
The independently evaluated quotient after dividing by
$2\widehat h(P)\prod_pc_p/\#E(\Q)_{\rm tors}^2$ is one, to
more than 69 digits, for $i=2,3,4,5,6,7,9,10$, and nine for
$i=11$. Thus the numerically predicted rational factors $q_i$ are
\[
 8,4,32,4,8,4,16,16,576,
\]
in that order. This initial numerical comparison alone proves neither
these exact factors nor the orders $1$ or $9$ of $\operatorname{Sha}$.
The exact factors are subsequently proved in C.6.R4 and C.6.R5:
the latter supplies an independent Gross--Zagier denominator bound
and rigorous intervals for all remaining ratios. This determines
the analytic BSD quotients exactly, while the orders of
$\operatorname{Sha}$ are proved here only in the three cases of
C.6.R4. In particular primitivity alone does not determine
$\#\operatorname{Sha}$ for the last curve.

\begin{proposition}[ARI C.6.R4; Known with derivation: three exact leading factors]\label{ari:C.6.R4}
For $i=3,5,7$, namely $j=287496,-3375,-884736$, one has
\[
 \operatorname{Sha}(E_i/\Q)=0,\qquad
 \frac{L'(E_i,1)}{\Omega_i}=4\widehat h_O(P_i).
\]
Thus the factors $q_3=q_5=q_7=4$ are exact, not numerical
rational reconstructions. For the smallest of these examples,
$E:y^2+y=x^3-38x+90$ and $P=(0,9)$, this gives the exact
reduction
\[
 N(P)=\exp\left(-\frac{L'(E,1)}{4\Omega_E}\right).
\]
\end{proposition}
\begin{proof}
We use \citet[Theorem 1.2 and Table 9]{miller2011}. In the
notation of its Definition 1.1, the theorem proves
$\mathrm{BSD}(E/\Q,p)$ for every prime $p$ when the curve has
analytic rank at most one, conductor less than $5000$, and
its isogeny class is outside the eleven exceptions in Table 9.
The exception conductors are
\[
546,570,858,870,1050,1230,1938,1950,2370,2550,3270.
\]
The three conductors here are $3136,3136,361$, and their
analytic ranks are one by C.6.R2. Hence every hypothesis is
verified. The theorem is stated for all curves in the isogeny
class, using isogeny invariance, so no optimal-curve assumption
is being silently imposed on our minimal equations.

In all three rows
$2\prod_pc_p/\#E(\Q)_{\rm tors}^2=4$, and primitivity is
proved in C.6.R1. Therefore the quotient
$L'(E_i,1)/(4\Omega_i\widehat h_O(P_i))$ is the positive
integer $\#\operatorname{Sha}(E_i/\Q)$. To isolate this integer
without relying on an unbounded decimal approximation to a
canonical height, we give an elementary coarse height certificate.

After completing the square, write the duplication map on the
$x$-line as $[F(X,Z):G(X,Z)]$, with relatively prime integral
quartics. The program
\path{code/C_height_coarse_certificate.py} supplies integer
homogeneous cubics $A_0,B_0,A_1,B_1$ and $D\ne0$ with
\[
 A_0F+B_0G=DZ^7,\qquad A_1F+B_1G=DX^7.
\]
Let $C$ be the larger of the coefficient one-norms of $F,G$
and of the two pairs $(A_j,B_j)$. For coprime integers $a,b$,
the common divisor of $F(a,b),G(a,b)$ divides $D$, by these
identities. Taking the identity corresponding to the larger
of $|a|,|b|$ proves, together with the elementary upper bound,
\[
 |h(x(2Q))-4h(x(Q))|\leq\log C.
\]
Telescoping consequently gives
\[
 \left|\widehat h_O(P)-\frac{h(x(2^6P))}{2\cdot4^6}\right|
 \leq\frac{\log C}{6\cdot4^6}.
\]
The exact constants in row order $3,5,7$ are
\[
 C=26610617537591808,\quad188716953600,\quad285931133.
\]
Six rational doublings, the displayed identities, and directed
Arb logarithms give the quotient enclosures
$1\pm0.003$, $1\pm0.003$, and $1\pm0.009$, respectively,
after including the central-derivative error bound from C.6.R2.
Each interval contains only the integer one. All quartics,
Bezout coefficients, rational iterates and enclosures are
retained in the program's output and its JSON certificate.
Finally C.6.3 gives $A=1$ in row seven, proving the exponential
formula. Its algebraicity remains the separate question C.7(c).
\end{proof}

\paragraph{The scope of Rubin's theorem.}
The applicable ordinary CM results concern a prime $p$ at which
$E$ has good ordinary reduction, equivalently $p$ splits in the CM
field. One precise opened formulation of Rubin's rank-one formula
is \citet[Theorem 1]{CrBDP2012}: for a curve over $\Q$ with CM
by the maximal order, an ordinary prime prime to the conductor,
and a simple central zero, an appropriate Katz $p$-adic $L$-value
is, up to a nonzero algebraic factor, the square of a formal
logarithm of a nontorsion rational point divided by a $p$-adic
period. Its input is analytic rank one, already verified here.
It is a $p$-adic identity; it does not identify the real exponential
$\exp\widehat h(P)$ as algebraic.

For the $p$-part of the real leading-term formula, a modern precise
version is \citet[Theorem A]{CrCastella2025}. In addition to analytic
rank one it requires full CM, abelian torsion field over the CM
field, $p\nmid6h_K$, splitting of $p$, conductor prime to $p$,
and that the different of $K$ divide the Hecke conductor exactly.
For the maximal odd-discriminant cases here, the latter follows
from the exponent two of the ramified prime in the rational
conductor. The nonmaximal cases of discriminant $-12$ and $-28$
can be passed through a $2$-isogeny at $p>3$. The discriminant
$-16$ case requires care: the direct application of this version
has a different-conductor hypothesis at $2$ that fails, and we
do not claim that this particular formulation proves its full
leading-term formula. Rubin's original ordinary main-conjecture
paper is \citep{CrRubin1991}; its bibliographic record was checked,
but its complete theorem text was not available in the opened
publisher copy, so no additional hypothesis-free claim is based
on it here.

In particular these ordinary-prime statements do not prove that
the $3$-primary part of $\operatorname{Sha}$ has order $9$ for the
last curve: $3$ is inert in $\Q(\sqrt{-163})$ and also excluded
from the stated $p>3$ hypotheses. C.6.R5 proves the corresponding
analytic quotient to be exactly $9$ by a separate denominator
argument. The identification with the actual Tate--Shafarevich
group remains a separate arithmetic question in that row.
